% SPDX-License-Identifier: Apache-2.0
% Build with lualatex (tools/build.sh --pdf sets document_build=lualatex):
% the theory text contains CJK characters, typeset via luatexja with an IPA font.
\documentclass[11pt,a4paper]{article}
\usepackage{fontspec}
\usepackage{luatexja-fontspec}
\setmainjfont{IPAMincho}
\setsansjfont{IPAGothic}
\usepackage{isabelle,isabellesym}
\usepackage{amssymb}
\usepackage{pdfsetup}

\begin{document}

\title{Buddhist and Comparative Asian Logic\\
Formalizations\\
Isabelle/HOL Technical Reference Manual\\[0.5ex]
\large\texttt{Buddhist\_Comparative\_Logic}}
\author{}
\date{}
\maketitle

\section*{About this manual}
This technical reference manual is the generated theory listing for the
\texttt{Buddhist\_Comparative\_Logic} Isabelle session in the Buddhist and
Comparative Asian Logic Formalizations software corpus. The session includes
generic logical infrastructure and formalizations across Buddhist and
comparative Asian traditions. The Heart Sutra is one detailed textual case
study: it uses a project-adopted 262-character received witness of the
Xuanzang-attributed Prajnaparamitahrdaya, not the 260-character Taisho body.
The reconstruction is parametric over classical, Belnap--Dunn FDE, and
Priest's FDE-plus-fifth-value semantics. This manual is not a paper or
preprint, and it is not a proof of any Buddhist doctrine; see
\texttt{docs/STATUS.md} for the wording rules and the clause-by-clause
coverage table.

\tableofcontents

% generated text of all theories
\input{session}

\nocite{*}
\bibliographystyle{abbrv}
\bibliography{root}

\end{document}
